\subsubsection{Caractère d'action et provenance de ses coefficients}
\label{sec:revision-planck}

La constante de Planck réduite, \(\hbar\), est l'objet physique qui
motive la réalisation de la section d'action angulaire. Sa position dans
la dépendance électronique doit être spécifiée à deux niveaux. La norme
déterminantale \(\Sigma_H\) détermine la décade de l'échelle. Le caractère
\(H_5\) et le retour régional déterminent ensuite les coordonnées des
sections d'action. Cette séparation permet d'exposer le résultat
nécessaire à l'électron sans anticiper les développements complets
d'autres secteurs de la Mécanique Dimensionnelle.

La provenance du caractère d'ordre cinq peut être explicitée au moyen
d'invariants déjà fixés dans le support. Dans la coordinatisation centrale, on considère
\[
 B_c=\begin{pmatrix}7&2\\2&7\end{pmatrix},
 \qquad \lambda_+=9,\quad\lambda_-=5.
\]
Le calendrier est de longueur \(12\), l'orbite du pas \(3\) est de
longueur \(4\), et le quotient de dénombrement pertinent est
\(104976/1944=54\). Les coefficients du caractère satisfont
\begin{equation}
 \frac9{16}=\frac{\lambda_+}{4^2},\quad
 \frac59=\frac{\lambda_-}{\lambda_+},\quad
 \frac7{48}=\frac{\operatorname{tr}(B_c)/2}{4\cdot12},\quad
 \frac1{54}=\frac{1944}{104976}.
 \label{eq:revision-h5-coeficientes}
\end{equation}
L'attribution de ces invariants aux degrés \(2,3,4,5\), avec
l'alternance des signes de \eqref{eq:el-h5}, appartient à la définition du
caractère analytique utilisé. L'expliciter est essentiel : l'ensemble des
cardinaux, considéré sans le caractère ni l'ordre de ses degrés,
admettrait d'autres expressions. Le contenu constructif réside dans la
composition du support, de l'orientation et de la règle analytique déclarée.

\subsubsection{Continuation régionale et équation de renouvellement}

Le sélecteur \(169\) détermine l'impulsion de degré six. Le prolongement
ultérieur s'organise selon le transfert \(D_A\) et une normalisation
unitaire de la droite déterminantale. Ces règles admettent une formulation
en séries formelles avant la substitution \(z=\alpha\), avec \(D_A\) déjà
construit. Écrivons
\[
 \mathscr C(z)=169z^6+R(z),\qquad R(z)\in z^7\mathbb R[[z]].
\]
La concaténation d'un prolongement augmente le degré d'une unité et
multiplie son poids par \(D_A\). Le premier prolongement strict a pour
poids \(D_A\). Son équation de renouvellement est donc
\begin{equation}
 R(z)=D_Az^7+D_AzR(z).
 \label{eq:revision-renovacion}
\end{equation}

\begin{proposition}[Unicité formelle de la continuation normalisée]
L'équation \eqref{eq:revision-renovacion} détermine une unique série et
\[
 \mathscr C(z)=169z^6+\frac{D_Az^7}{1-D_Az}
             =169z^6+\sum_{n=7}^{\infty}D_A^{n-6}z^n.
\]
La réalisation est analytique pour \(|D_Az|<1\).
\end{proposition}
\begin{proof}
L'élément \(1-D_Az\) a pour terme constant un et est inversible
dans l'anneau des séries formelles. Résoudre l'équation produit la série
affichée. De manière équivalente, le coefficient de degré sept vaut \(D_A\)
et chaque coefficient ultérieur vaut \(D_A\) fois le précédent. La série
géométrique converge absolument dans le disque indiqué.
\end{proof}

La normalisation du premier prolongement est une donnée opératoire
indispensable. Si l'on ne conservait que l'impulsion de degré six et
le rapport de concaténation, les séries
\[
 169z^6+\lambda\frac{D_Az^7}{1-D_Az}
\]
partageraient encore ces deux données pour tout \(\lambda\).
La règle unitaire fixe \(\lambda=1\) avant l'évaluation.
La contribution de chaque règle à l'unicité de l'opérateur de lecture
est ainsi localisée, et l'expression rationnelle conserve la continuation complète.

\subsubsection{Constante de Planck réduite et retour de la section d'action}

Dans la coordinatisation dimensionnelle de \eqref{eq:el-secciones-accion}, la section
antérieure au retour est la construction du modèle destinée à la comparaison
avec la constante de Planck réduite :
\[
 \hbar_{\mathrm{pre}}=H_5\,10^{-34}\mathcal U_S.
\]
La section postérieure conserve également la compensation régionale
\(90\mathscr C(\alpha)/\pi\). Toutes deux appartiennent à la même droite
d'action et leur quotient \(R_{\mathrm{act}}\) mesure l'effet multiplicatif
du retour. Le passage à \(h\) par \(2\pi\) utilise la relation
entre la paramétrisation angulaire et un tour complet.

Cette construction intervient effectivement dans
\(R_{\mathrm{act}}^{80-54\alpha+6\Delta_4}\). L'électron conserve
sa coordonnée centrale et sa trajectoire ; l'échelle chronologique ne
remplace pas cette structure par un quantum indifférencié. On ne
multiplie pas non plus à nouveau par \(10^{-34}\) une masse déjà exprimée en MeV :
la décade appartient à la section d'action, se simplifie dans son quotient
sans dimension et ne réapparaît que là où agit la réalisation
dimensionnelle correspondante.

La comparaison numérique de \(H_5\), de la section retournée et de
\(h/(2\pi)\) est présentée dans la section~\ref{sec:revision-planck-contraste}.
On distingue ainsi la dénomination de l'objet physique, l'équation que le
modèle lui attribue et la précision avec laquelle cette équation le reproduit.
